\documentclass[UTF8,a4paper,12pt]{ctexart}
\usepackage{geometry}
\usepackage{hyperref}
\usepackage{longtable}
\usepackage{xcolor}
\geometry{left=2.4cm,right=2.4cm,top=2.4cm,bottom=2.4cm}
\hypersetup{colorlinks=true,linkcolor=blue,urlcolor=blue}
\setlength{\parindent}{2em}
\setlength{\parskip}{0.35em}

\title{\textbf{MoonRouteKit 项目申报书}\\MoonBit 路由模式编译与请求匹配内核}
\author{python123}
\date{2026 年 6 月}

\begin{document}
\maketitle
\tableofcontents
\newpage

\section{项目名称}
MoonRouteKit：MoonBit 路由模式编译与请求匹配内核。

\section{仓库链接}
\begin{itemize}
  \item GitLink 仓库：\url{https://gitlink.org.cn/python123/moonroutekit}
  \item GitHub 仓库：\url{https://github.com/python123-ops/moonroutekit}
\end{itemize}

\section{项目简介}
MoonRouteKit 是一个面向 MoonBit 生态的轻量级 URL 路由核心库，提供路径模式编译、HTTP 方法匹配、参数提取、命名路由反向生成和路由冲突诊断能力。项目定位在 Web 框架底层的可复用内核，不绑定具体服务器、运行时或响应模型，适合 MoonBit 服务端示例、边缘计算实验、Wasm 服务原型和教学场景复用。

\section{项目方向与适用场景}
MoonBit 生态正在扩展到浏览器、Wasm、服务端实验和工程化教学等场景。很多示例项目需要把请求路径映射到处理逻辑，但直接实现完整 Web 框架成本较高。MoonRouteKit 提供一个小而稳定的路由匹配核心，使开发者可以在不同运行时中复用同一套路由规则。

适用场景包括：
\begin{itemize}
  \item MoonBit HTTP 服务示例的请求分发。
  \item 边缘函数和 Wasm 服务中的路径匹配。
  \item 网关原型、课程实验和框架教学中的路由内核。
  \item 需要命名路由反向生成的轻量级服务组件。
\end{itemize}

\section{拟实现的核心功能}
\begin{longtable}{p{4cm}p{10cm}}
\textbf{功能模块} & \textbf{说明} \\
路径模式编译 & 支持根路径、静态路径、参数片段和最终通配片段。 \\
请求匹配 & 支持 HTTP 方法过滤和确定性优先级选择。 \\
参数提取 & 将匹配得到的动态片段输出为稳定键值对。 \\
反向生成 & 根据命名路由和参数生成 URL。 \\
冲突诊断 & 检测重复路由和通配路由遮蔽问题。 \\
CLI 示例 & 展示匹配、参数提取、反向生成和诊断输出。 \\
\end{longtable}

\section{技术路线}
MoonRouteKit 将路径模式编译为片段数组，匹配时先过滤方法，再逐段比较请求路径。静态片段获得最高优先级，参数片段次之，通配片段最低。该规则保证匹配结果可预测，也便于测试和审查。

反向生成阶段通过命名路由定位编译后的路径模式，将参数片段替换为调用方传入的值。若缺少必需参数，函数返回结构化错误，避免产生不完整 URL。

诊断阶段扫描路由表，发现重复路径和通配遮蔽。诊断结果可在 CLI、测试或上层框架初始化阶段展示。

\section{原创说明}
本项目为原创实现，不移植既有开源项目源码。项目选题聚焦 MoonBit Web 与边缘服务实验中的路由内核，区别于数据表处理、日志查询、项目元数据检查、构建证据记录、数据流分析等方向。实现不引入外部依赖，核心算法、测试样例和文档均围绕 MoonBit 生态的轻量级服务组件需求设计。

\section{阶段成果}
当前仓库包含 MoonBit 核心库代码、黑盒测试用例、CLI 演示程序、设计文档、验收清单和申报书材料。后续可继续补充基准测试、更多路径规范策略和与 MoonBit HTTP 示例项目的集成演示。

\end{document}
